\section{Conclusions and Future Work}
\label{sec:conclusions}

GeoSoil-Cuba demonstrates an open-source geospatial
workflow for preparing historical soil cartography,
integrating soil and vegetation attributes, and
examining Sentinel-2 NDVI observations within an
agricultural landscape in Jaruco, Mayabeque, Cuba.
The analysis combines vector-data validation,
coordinate reference system harmonization,
spatial overlay, satellite-image processing,
and temporal vegetation-index analysis within
a documented GIS and Python environment.

Within the 178.47~km$^2$ study area, the soil
coverage assessment reported approximately
162.85~km$^2$ of mapped soil units. Accounting
for mapped water surfaces reduced the remaining
non-water soil-data gaps to approximately
4.20~km$^2$, equivalent to 2.35\% of the AOI.
Topology assessment identified 90 overlapping
polygon pairs with a cumulative pairwise overlap
area of approximately 0.054919~km$^2$.
These findings highlight the importance of
evaluating spatial coverage, geometric consistency,
and area-accounting procedures before interpreting
historical soil-map information.

Satellite processing produced 32 regional NDVI
rasters, of which 15 acquisitions satisfied
the regional cloud-coverage criterion.
Within the selected 16.22~km$^2$ agricultural
region, mean NDVI declined between November 2025
and March 2026, followed by fluctuations and
an increase by June 2026. A complementary
analysis of three locally verified potato fields
revealed different temporal NDVI trajectories,
illustrating the importance of distinguishing
field-level observations from regional aggregates.

The findings characterize spatial and temporal
patterns in the available datasets but do not
establish crop phenological stages, agricultural
productivity, or causal relationships between
soil classes and vegetation conditions.
Interpretation remains constrained by uncertainties
in historical mapping, unresolved topology and
area-accounting discrepancies, irregular satellite
sampling, residual cloud contamination, and
incomplete field and radiometric metadata
verification.

Future development should prioritize reconciling
soil and vegetation area calculations, reviewing
overlapping soil polygons and remaining non-water
gaps, verifying radiometric parameters for
individual satellite products, and implementing
explicit pixel-level cloud and shadow masking.
Additional improvements should include accurate
agricultural field boundaries, planting and
harvesting records, and ground observations
corresponding to satellite acquisition dates.
These developments would provide a stronger
basis for evaluating vegetation variability
across mapped soil units and agricultural
management conditions.

The resulting framework provides a foundation
for further geospatial research and potential
application in other regions where historical
soil cartography, ancillary land-cover data,
and contemporary satellite observations must
be integrated while preserving data provenance
and documenting uncertainty.